\begingroup
\clearpage
\let\clearpage\relax
\vspace*{-40pt}
\chapter[BACKGROUND AND MOTIVATION]{BACKGROUND AND MOTIVATION}\label{sec.background}
\endgroup

This chapter motivates the study undertaken in this dissertation and covers necessary background.
Section~\ref{sec.ehr} introduces the role of electronic health records
in comorbidity analysis and mortality risk assessment.
Section~\ref{sec.mcp} reviews classical and modern approaches to the
maximum clique problem. Section~\ref{sec.fp} surveys fractional
programming and its integer and combinatorial extensions.
Section~\ref{sec.contributions} summarizes the contributions of this
dissertation.

\section{Electronic Health Records in Comorbidity and Mortality
Analysis}\label{sec.ehr}

Electronic health records (EHRs) capture comprehensive longitudinal
patient information, including demographics, diagnoses, medications,
laboratory results, and procedures. These rich, multimodal datasets
support large-scale cohort studies, knowledge discovery, and the
development of longitudinal phenotypic profiles that inform clinical
decision-making and personalized medicine~\citep{Jensen2012}. They
have become indispensable for healthcare analytics, enabling advanced
data mining techniques such as information extraction, representation
learning, outcome prediction, phenotyping, de-identification, disease
prediction, comorbidity analysis, and patient
stratification~\citep{Yadav2018,Sarwar2023}.

The widespread adoption of EHR systems, accelerated by initiatives
such as the U.S. Health Information Technology for Economic and
Clinical Health Act of 2009, has produced vast repositories of
structured and unstructured data. By 2015, 84\% of hospitals and 87\%
of office-based physicians had implemented certified EHR
systems~\citep{Shickel2018}. This expansion has extended EHR utility
beyond primary functions such as billing and documentation to
secondary applications in clinical informatics, including the
extraction of medical concepts from clinical notes~\citep{Sarwar2023},
learning robust patient representations for predictive
modeling~\citep{Shickel2018}, forecasting outcomes such as hospital
readmission or heart failure~\citep{Yadav2018}, and automating
phenotyping for research cohort selection~\citep{Jensen2012}.

\textit{Mortality rate}, the fraction of deceased patients among those
with a particular disease or group of diseases, is a metric of broad
interest in medicine, epidemiology, and public health. The co-occurrence of multiple diseases in a single patient is often described using the related concepts of \textit{comorbidity} and \textit{multimorbidity}. Comorbidity refers to the presence of one or more additional conditions alongside a primary or index disease, whereas multimorbidity refers to the coexistence of two or more chronic conditions in an individual without designating any condition as primary. Comorbidity assumes particular significance in the context of mortality, as the interaction between two or more diseases can produce health outcomes that differ from those of each
disease in isolation~\citep{FEINSTEIN1970455}. Among the many
analytical opportunities offered by EHRs, comorbidity analysis stands
out as especially valuable for understanding mortality risk since
co-occurring diseases often interact in ways that elevate mortality
well beyond the simple additive effects of individual
conditions~\citep{FEINSTEIN1970455}. As \cite{Redelmeier1998}
highlight, a dominant illness can overshadow concurrent conditions and
increase the risk of clinical neglect, a challenge that data-driven
EHR analytics can help address.

Identifying multimorbidities in the form of disease clusters with high
mortality rates yields benefits across epidemiology, health policy,
healthcare management, and public health planning. Knowledge of such
clusters enables healthcare providers to anticipate disease onset early
and intervene proactively, improving the likelihood of successful
treatment and reducing mortality. It also supports more efficient
resource allocation, as health system administrators can direct
resources toward conditions most likely to lead to death. Furthermore,
high-mortality cluster identification can catalyze research into
underlying causes and potential treatments, and public health officials
can use this information to design targeted interventions for
high-risk populations.

\cite{Gijsen2001} conducted a comprehensive review of the causes and
consequences of comorbidity, finding that comorbidity was consistently
associated with increased mortality. Prior research indicates that
patients diagnosed with conditions such as myocardial infarction,
diabetes, chronic obstructive pulmonary disease, chronic kidney
disease, dementia, depression, hip fracture, stroke, colorectal
cancer, and lung cancer are at elevated risk of developing heart
failure~\citep{Ahluwalia2012}, a particularly notable link given that
heart disease and cancer account for the majority of reported deaths
in the United States~\citep{Xu2016}. \cite{Braunstein2003} further
show that elderly patients with chronic heart failure often succumb to
non-cardiac conditions, suggesting that identifying and addressing
these conditions can improve outcomes. \cite{CORRAINI2018242}
investigated the impact of comorbidity on post-stroke mortality in a
population-based cohort study, finding that conditions such as cancer,
advanced renal disease, or liver disease significantly increased
one-year mortality following a stroke.

Analytical approaches to mortality and comorbidity have historically
relied on statistical methods, with EHR data now enabling more
granular, longitudinal investigations. Early work by \cite{lu1996}
introduced a mixed Poisson regression model for age-specific and
age-standardized mortality rates, using a marginal quasi-likelihood
approach within an empirical Bayes framework to improve practical
analysis, particularly for diseases such as lung cancer.
\cite{LU2021113583} modeled comorbidities as temporal disease networks
to visualize comorbidity progression and identify critical transition
points using clustering techniques. More recent studies have
incorporated comorbidities to enhance predictive performance. For
example, \cite{ZOLBANIN2015150} demonstrated that including
comorbidities improves survival prediction models for cancer patients,
while gender-focused analyses have explored related
factors~\citep{Short2013S93}. Many investigations target comorbidities
linked to specific conditions, including diabetes, obesity, and
cancer~\citep{Holguin2005, Marrie2015, Leontiadis2013}.
\cite{vanGestel2013} found that comorbidity burden and advanced age
were associated with early postoperative mortality in gastrointestinal
cancer patients, with cardiovascular diseases exerting a particularly
strong effect on 30-day mortality. \cite{Koskinen2022} employed a
data-driven approach to analyze comorbidity patterns across a range of
common diseases, identifying disease-specific patient subgroups with
distinctive diagnosis patterns, survival functions, and laboratory
correlates, underscoring the heterogeneity within disease populations
and the need for personalized risk assessment.

Recent EHR-based studies have employed association rule mining (e.g. the a priori algorithm) and $k$-means clustering to uncover
frequent disease co-occurrences and support early detection and
tailored treatment~\citep{HariEtAl2024INFORMStutorial,Parisa2025}.
Optimization frameworks offer additional promise by formulating complex
analytical tasks as constrained problems. \cite{Zhong2022} modeled
relational-sequential patient clustering with custom distance measures
that account for demographic hierarchies and temporal diagnosis
similarity, producing clinically interpretable clusters. Similarly,
\cite{VaghfiMohebbi2025lethalclique} defined the maximum mortality
rate clique problem and applied enumeration to identify high-mortality
disease cliques in comorbidity graphs. 

The following section reviews the maximum clique problem and related solution methods that underpin this line of work.

\section{The Maximum Clique Problem}\label{sec.mcp}

A \emph{clique} in an undirected graph $G=(V,E)$ is a subset of
vertices in which every pair is connected by an edge. The \emph{maximum
clique problem} (MCP) asks for the largest such subset, i.e., the
clique of maximum cardinality. Closely related are the \emph{maximum
independent set problem} (MIS), which seeks the largest subset of
mutually non-adjacent vertices (also known as the stability number),
and the \emph{minimum vertex cover problem}, which seeks the smallest
set of vertices incident to every edge. These three problems are
equivalent in the following sense: a set $S$ is a maximum clique in
$G$ if and only if $S$ is a maximum independent set in the complement
graph $\bar{G}$, and if and only if $V \setminus S$ is a minimum
vertex cover in $\bar{G}$. This equivalence has enabled substantial
transfer of algorithmic ideas and bounding techniques across
all three problems~\citep{Bomze99survey}.

All three problems are NP-hard on general graphs~\citep{GareyJohnson79},
which has motivated a rich body of research on both exact and
approximate solution methods. The maximum clique problem admits several
natural integer programming formulations. The most basic is the
edge formulation, which maximizes the number of selected
vertices subject to pairwise incompatibility constraints derived from
the complement graph:
\begin{subequations}\label{form.eclique}
    \begin{align}
        \max &\ \sum_{i=1}^n x_i \label{ec.obj}\\
        \text{s.t.}\quad
        & x_i + x_j \le 1 && \forall \{i,j\} \in E(\bar{G})
            \label{ec.constraint}\\
        & x_i \in \{0,1\} && \forall i \in V(G).
    \end{align}
\end{subequations}
In formulation~\eqref{form.eclique}, $x_i=1$ means vertex $i$ is in the clique, and the constraint forbids selecting both endpoints of any non-edge, so any feasible binary solution corresponds to a clique. 
% \comm{You defined a clique a subset of nodes; so a binary vector cannot be one.}

More advanced mixed-integer formulations---including
clique-inequality-based, independent-set-augmented, and lifted
variants---have been developed to strengthen the LP relaxation and
improve branch-and-bound performance; a detailed comparison of their
strength and computational behavior appears in~\citep{Naderi2022}.
Earlier works explored quadratic programming representations and
exploited the Motzkin--Straus theorem, which links the maximum clique
size to the optimal value of a continuous quadratic program over the
simplex~\citep{MotzkinStraus65, pardalos1987constrained, pardalos1994}.

Comprehensive surveys of algorithmic progress on the maximum clique
problem include~\cite{PardalosXue1994}, \cite{Bomze99survey},
\cite{WU2015693}, and the more recent overview by~\cite{marino2024}.
Drawing from these, we organize solution methods into three principal
categories: enumerative algorithms, exact implicit-search approaches
(including branch-and-bound, cutting-plane, and reformulation
techniques), and heuristic methods.
\paragraph{Enumerative algorithms.}
Enumerative algorithms form the classical foundation for exact
solutions, systematically exploring the search space to list all
maximal cliques---or identify the maximum one---without repetition or
omission. These methods typically rely on backtracking with
pruning to manage exponential growth.

The earliest explicit enumeration of all cliques appeared in
\cite{HararyRoss57}, which introduced a basic listing procedure based
on adjacency relations. In the 1960s, these ideas were extended to
clustering and information-retrieval applications: \cite{Bonner64} and \cite{BednarekTaulbee66} developed among the first procedures to enumerate all maximal cliques by repeatedly extending partial cliques against the adjacency structure, rather than locating a single clique.

A significant advance came in the early 1970s with efficient
backtracking schemes. \cite{Akkoyunlu73} described an early recursive
procedure that pruned invalid branches using adjacency checks, while
the now-classic \cite{BronKerbosch73} algorithm introduced a
pivot-based strategy that partitions the search space by selecting a
pivot vertex, thereby eliminating redundant recursive calls and
achieving polynomial space complexity. This pivot mechanism proved both
practical and highly influential.

Subsequent refinements in the 1970s built on the Bron--Kerbosch
framework. Decomposition methods confined maximal cliques to individual
subgraphs for modular or parallel computation~\citep{meeusen1975};
lexicographic enumeration of maximal independent sets used careful
depth-first candidate management for output-sensitive
performance~\citep{johnson1975}; systematic variant comparisons
confirmed the original pivot approach as highly
competitive~\citep{johnston1976}; and a forward-checking strategy that tests candidate vertices for extendability {before} recursing pruned the search earlier than pivoting alone, improving on Bron--Kerbosch on sparse graphs~\citep{gerhards1979}.

Research in the 1980s and later decades shifted toward tighter pruning
and improved ordering. Depth-first algorithms with aggressive
lexicographic ordering and candidate elimination outperformed earlier
methods~\citep{loukakis1981}, while a listing procedure of \cite{chiba1985} organized the enumeration around the graph's sparse structure to remain efficient on sparse graphs. A major contribution came from
\cite{TOMITA2006}, who developed enhanced Bron--Kerbosch variants whose pivot selection reduces the number of recursive branches, yielding substantially faster practical performance and what is widely regarded as one of the strongest enumeration schemes for general graphs. Later work refined this lineage by ordering the recursion along a degeneracy ordering of the vertices—processing low-degree vertices first—to bound the work per branch on sparse graphs~\citep{CAZALS2008TCS, Eppstein13}; further practical gains came from combining degeneracy ordering, parallelization, and tighter bounding on sparse real-world instances~\citep{san2018efficiently}.

\paragraph{Exact implicit-search methods.}
While enumerative methods provide complete listings, many applications require only the single largest clique and benefit from implicit enumeration techniques that
avoid exhaustive listing. These approaches, most commonly realized
through branch-and-bound (B\&B) frameworks, systematically construct
partial solutions by branching on vertex inclusion or exclusion while
using bounding functions, often graph coloring or related upper-bound
estimators, to prune branches that cannot improve upon the current
best solution. This combination guarantees optimality while exploring
a dramatically smaller portion of the search space than naive
enumeration~\citep{PardalosXue1994}.

Early implicit enumeration efforts in the 1970s introduced basic
backtracking with rudimentary bounding for both MCP and
MIS~\citep{desler1970, tarjan1972, HouckVemuganti77}, subsequently
strengthened by more effective lower bounds on the independence number
derived from clique cover arguments~\citep{TarjanTrojanowski77}. In
the 1980s, graph-theoretic properties were exploited more aggressively:
one approach extracted a large induced triangulated (chordal) subgraph, within which a maximum clique can be found efficiently, and used it to bound and prune the surrounding search~\citep{BalasYu86}; another used the
number of incident triangles as a branching criterion to prioritize
highly connected vertices~\citep{gendreau1993solving}.
\cite{CarraghanPardalos90} developed a depth-first algorithm that prunes a branch once the current clique plus the remaining candidate vertices cannot exceed the incumbent; its degree-based ordering and low memory use made it a long-standing benchmark on sparse graphs.

Contemporary exact algorithms increasingly rely on upper bounds
obtained by coloring induced subgraphs of candidate sets. The Cliquer algorithm of \cite{Ostergard02} processes vertices in a fixed order and reuses a bound computed for previously considered vertices to prune branches that cannot improve the incumbent. Similarly, \cite{TomitaSeki2003} uses
coloring of the remaining candidate set both to estimate an upper bound
via the chromatic number and to provide a branching heuristic by partitioning vertices into color classes. This coloring-based bounding strategy has become central to state-of-the-art B\&B solvers.

Further refinements have strengthened coloring integration through
strategies such as DSATUR-style greedy coloring with degeneracy
ordering, bit-parallel implementations, and dynamic recoloring,
producing tighter and faster upper
bounds~\citep{bourjolly1996, BOURJOLLY199721, Kumlander2005,
TomitaKameda2007, Konc2007AnIB, Tomita2010, SANSEGUNDO2011571}.
Complementary work incorporates MaxSAT-based reasoning to detect
infeasibility and accelerate optimality proofs on hard
instances~\citep{LiQuan2010, WU2015693, Maslov2014}. A parallel line
of research has advanced branch-and-cut methods for the equivalent
maximum independent set formulation, using edge-projection inequalities
to strengthen LP relaxations~\citep{Mannino1996, Rossi2001,
Rebennack2011}, with comprehensive surveys available
in~\cite{Rebennack2012}.

\paragraph{Heuristic methods.}
For large-scale graphs where exact methods become computationally
prohibitive, heuristic methods provide a practical alternative,
forgoing provable optimality in favor of high-quality solutions within
acceptable running times. The simplest approaches are sequential greedy
constructions that build maximal cliques incrementally by repeatedly
selecting the vertex that optimizes a local criterion (e.g., degree in
the remaining subgraph or number of common neighbors). Influential
early examples include the greedy procedures of \cite{johnson1975} and
vertex-ordering-based methods by \cite{tomita1988}, which remain fast
baselines owing to their low implementation overhead.

Local search heuristics extend this idea by starting from an initial
clique and iteratively exploring its neighborhood---through vertex
addition, removal, or swap operations---until no improving move
exists~\citep{foe1994}. These methods frequently outperform pure greedy
approaches by permitting small escapes from poor local decisions,
though they can still be trapped in suboptimal local maxima.

Metaheuristic frameworks address this limitation by introducing
controlled diversification and intensification. Tabu search employs
short-term memory structures (tabu lists) to prevent revisiting recent
solutions, thereby avoiding cycling and promoting broader
exploration~\citep{FridenA89, gendreau1993solving, Battiti99,
soriano1996}. Simulated annealing probabilistically accepts worsening
moves according to a decreasing temperature schedule, enabling escape
from shallow local optima early in the
search~\citep{AartsKorst89, jerrum92, homer94}. Genetic algorithms
evolve a population of candidate cliques through fitness-based
selection, crossover, and mutation, often delivering strong results on
structured or moderately sized
graphs~\citep{back1994, carter1993, foster1995, Guturu1994CliqueFG,
marchiori98}.

Neural network-based methods offer a further perspective, particularly
for exploiting graph structure. Early connectionist approaches used
Hopfield-style recurrent networks to iteratively relax toward maximal
cliques~\citep{ramanujam1988, Jagota95, jagota1992, AartsKorst89,
bertoni1997, FUNABIKI1992340, shrivastava1990, wu1994NewMT}. More
recently, graph neural networks (GNNs) have emerged as powerful tools:
by aggregating information from neighboring nodes, GNNs learn
expressive representations of vertex roles and subgraph patterns
directly from graph-structured data. Recent work has successfully
applied GNNs to predict high-quality cliques or guide search
procedures, with promising generalization across diverse graph
families~\citep{Karalias2020, min2022}.

A complementary approximation strategy operates at the subgraph level,
first extracting a smaller, more tractable subgraph (via degeneracy
ordering, community detection, or random sampling) and then solving
MCP exactly or heuristically within it~\citep{xue1994}. When the
selected subgraph preserves the key dense regions, this approach can
yield strong solutions at significantly reduced computational cost.

The algorithmic landscape of the maximum clique problem, spanning
enumeration, exact implicit search, and scalable heuristics, provides
powerful tools for identifying dense substructures in graphs. 

In the next section, we extend our perspective to ratio-based objectives by reviewing fractional programming and emphasizing combinatorial variants.

\section{Fractional Programming}\label{sec.fp}

Fractional programming concerns the optimization of ratios of functions
and has important applications in operations research, economics, and
engineering. The field traces its origins to
\cite{vonNeumann1937}'s equilibrium model for an expanding economy.
Systematic development began with the seminal work of
\cite{charnes1962}, who demonstrated that linear fractional programs
can be transformed into equivalent linear programs via a simple
variable substitution. Subsequent research situated fractional
programming within the broader framework of generalized
convexity~\citep{martos1975} and global optimization, recognizing that
ratio objectives often yield nonconvex problems with multiple local
optima. Comprehensive treatments of single-ratio fractional programs
appear in~\cite{schaible1978analyse, schaible1981generalized,
schaible1983fractional, craven1988fractional, stancu1997fractional},
while broader surveys and collections are available
in~\cite{FloudasPardalos02, schaible2002, schaibleshi2004, Frenk2005,
cambini2009generalized} and related
proceedings~\citep{cambini1990generalized, crouzeix1998generalized,
singh1989continuous, schaible1995fractional}. A focused
review of fractional approaches in location problems is given
in~\cite{barros1998discrete}.

Let $f$, $g$, and $h_k$ ($k=1,\dots,m$) be functions defined over a
set $D \subseteq \mathbb{R}^n$, with $g(x) > 0$ on $D$. Define the
ratio $q(x) = f(x)/g(x)$ and the feasible set
$S = \{x \in D : h_k(x) \le 0, \; k=1,\dots,m\}$.
The \emph{single-ratio fractional program} is:
\begin{align}
    \label{form.P}
    \sup \{ q(x) : x \in S \}.
\end{align}
When $f$, $g$, and the constraints are affine and
$D = \mathbb{R}_+^n$, this specializes to the
\emph{linear fractional program}:
\begin{align}
    \label{form.LCFP}
    \sup \left\{ \frac{c^T x + \alpha}{d^T x + \beta} :
        Ax \le b, \; x \ge 0 \right\},
\end{align}
with $c,d \in \mathbb{R}^n$ and $\alpha,\beta \in \mathbb{R}$.
Setting $d = \mathbf{1}$ and $\beta = 1$ yields the special case
known as \emph{uniform fractional optimization}.

Single-ratio fractional programs arise across a wide range of
application domains. In resource allocation and economic
decision-making, they are used to maximize profit-to-capital or
profit-to-revenue ratios~\citep{barros1998discrete, barros1997fractional,
mjelde1983methods}, expected return per unit of risk in portfolio
selection~\citep{schaible1995fractional}, output-to-input ratios in
data envelopment analysis~\citep{charnes1978measuring}, and labor
efficiency in warehouse operations~\citep{hoskins1984optimal}.
Inventory and maintenance models frequently minimize cost per unit
time~\citep{frenk1997marginal, barros1997optimizing}, cutting-stock
problems seek to minimize the ratio of wasted to utilized
material~\citep{GilmoreGomory1963}, and signal processing applications
maximize signal-to-noise ratios~\citep{falk1969maximization}.

In network flows, scheduling, and graph-theoretic optimization,
fractional programs appear in the maximum profit-to-time ratio cycle
problem and its dual~\citep{AhujaBook, Dantzig1967, Fox1969,
Lawler1967a, megiddo1979combinatorial, Megiddo1983}, the
minimum mean cycle problem~\citep{AhujaBook, Karp1978, Lawler1967a},
maximum mean-weight cuts~\citep{Radzik1993, Iwano1994, Ervolina1993},
fractional variants of the assignment
problem~\citep{shigeno1995algorithm}, and fractional binary knapsack
generalizations~\citep{Ishii1976, Hashizume1987}. Additional
applications span retail assortment planning, set covering,
biclustering, element dispersion, stochastic service systems, and the
generation of alternative solutions in binary linear
programs~\citep{borrero2017survey, radzik1998fractional}.

\paragraph{Solution methods for single-ratio programs.}
Several established techniques solve single-ratio fractional programs.
The classical approach of \cite{charnes1962} transforms
problem~\eqref{form.P} via the substitution $y = x / g(x)$ and
$t = 1 / g(x)$, yielding:
\begin{align}
\label{form.Peq}
    \sup \left\{ t f\left( \frac{y}{t} \right) :
        (y,t) \in \tilde{S} \right\},
\end{align}
where $\tilde{S}$ denotes the transformed feasible set and $t > 0$.
The resulting optimization problem is concave (linear if the original problem is linear fractional) and can therefore be solved efficiently using standard convex or linear programming solvers. If $(\bar{y}, \bar{t})$ is optimal for~\eqref{form.Peq}, then $\bar{x} = \bar{y} / \bar{t}$ solves~\eqref{form.P}. For linear fractional programs, this
transformation yields an equivalent LP solvable by the simplex method.

A second strategy exploits duality. Using the same substitution, the
Lagrangian dual of~\eqref{form.Peq} takes the form:
\begin{align}
\label{form.D}
     \inf \left\{ \sup_{x \in C} \frac{f(x) - u^T h(x)}{g(x)} :
         u \ge 0 \right\},
\end{align}
where $C$ is the original domain adjusted for constraints. Various
dual formulations have been proposed~\citep{schaible1978analyse,
schaible1983fractional}, with iterative bounding techniques developed
for concave fractional programs~\citep{schaible1976fractional}. Beyond
algorithmic utility, duality supports sensitivity analysis: the dual
reveals how perturbations in right-hand-side values affect the optimal
ratio.

The most widely used practical method parametrizes the problem with a
scalar $\lambda \in \mathbb{R}$:
\begin{align}
\label{form.Plam}
    \max \{ f(x) - \lambda g(x) : x \in S \}.
\end{align}
Under mild conditions (compact $S$, continuous $f$ and $g$), the
optimal value function $F(\lambda)$ is strictly decreasing and convex,
so $F(\lambda) = 0$ has a unique root $\lambda^*$. Solving
\eqref{form.Plam} at $\lambda^*$ yields an optimal solution to the
original ratio problem. The root can be located via bisection,
Newton's method, or modified interval methods~\citep{ibaraki1983parametric,
schaible1983fractional}. The Newton-based approach is Dinkelbach's algorithm~\citep{dinkelbach1967nonlinear}, which updates $\lambda$ to the ratio attained by the current solution and converges rapidly; variants with improved convergence followed~\citep{schaible1976fractional, pardalos1991global}.

Additional specialized methods include interior-point approaches for
linear fractional problems. Early work showed that Karmarkar's
projective algorithm can be interpreted as solving a linear fractional
program over the simplex~\citep{anstreicher1986monotonic}. Later
contributions adapted conjugate-gradient projection
schemes~\citep{TANTAWY2008969} and network flow techniques for certain
unconstrained binary fractional problems~\citep{picard1982network}.

\paragraph{Integer and combinatorial extensions.}
Many practical applications of fractional programming require
integer, typically binary, variables, particularly in economic
modeling, resource allocation, and combinatorial
optimization~\citep{barros1998discrete}. Comprehensive surveys of
integer fractional programming appear in~\cite{barros1998discrete} and,
more recently, in the context of hyperbolic binary
programs~\citep{borrero2017survey}.

A general \emph{fractional combinatorial optimization problem} (FCOP)
takes the form:
\begin{align}
    \label{form.FCOP}
    \max_{x \in \mathcal{X}} \frac{f(x)}{g(x)},
\end{align}
where $\mathcal{X} \subseteq \{0,1\}^n$ is a discrete feasible set
defined by combinatorial constraints, and
$f,g : \mathcal{X} \to \mathbb{R}$ with $g(x) > 0$ for all
$x \in \mathcal{X}$. When $f$ and $g$ are affine, this is known as a
\emph{single-ratio fractional binary program}~\citep{boros2002pseudo,
tawarmalani2002global, wu1997note}, and can often be reformulated as a
mixed-integer linear program (MILP) using linearization techniques or
auxiliary variables~\citep{LI1994590, tawarmalani2002global,
BORRERO2016479}.

FCOPs appear across diverse domains. A classical example is the minimum-ratio spanning-tree problem, which seeks a spanning tree minimizing the
ratio of total cost to total weight and motivated Megiddo's parametric-search technique, which transfers efficient algorithms for the linear-objective version of a combinatorial problem to its ratio-objective version~\citep{Chandrasekaran1977, megiddo1979combinatorial, Megiddo1983, skiscim2004complexity, ursulenko2013global}. In graph
theory, the \emph{maximum ratio clique problem} (MRCP) generalizes MCP~\citep{Sethuraman2015ratioclq}:
given vertex weights $a_i \ge 0$ (benefit) and $b_i > 0$ (cost), MRCP
finds a clique $C$ maximizing $\sum_{i\in C} a_i / \sum_{i\in C}
b_i$. When all weights are unitary,
MRCP reduces to the standard maximum clique problem. Specialized
approaches for MRCP include difference-of-convex (DC) programming with
valid inequalities to strengthen MILP relaxations~\citep{moeini2015}
and multi-start variable neighborhood search
heuristics~\citep{GOEKE2017283}.

The dominant generic approach for FCOPs is the parametric method,
which introduces a scalar $\lambda$ and solves:
\begin{align}
    \label{form.PFCOP}
    \max_{x \in \mathcal{X}} F(x) = f(x) - \lambda g(x).
\end{align}
Let $F(\lambda)$ denote its optimal value. Under standard assumptions,
$F(\lambda)$ is continuous, convex, piecewise linear, and strictly
decreasing, guaranteeing a unique root $\lambda^*$ equal to the
optimal ratio of the original FCOP. Because $\mathcal{X}$ is finite,
$F$ has at most $|\mathcal{X}|$ linear segments, enabling finite
convergence of root-finding procedures~\citep{schaible1983fractional,
Hashizume1987, grunspan1973hyperbolic, robillard1971hyperbolic}.

Two main strategies locate $\lambda^*$: binary search, which maintains
a bracketing interval and evaluates the sign of $F$ at the
midpoint~\citep{Radzik2009}; and Dinkelbach's
algorithm~\citep{dinkelbach1967nonlinear}, which in each iteration
solves~\eqref{form.PFCOP} at the current estimate $\lambda_i$ to
obtain $x^{(i)}$, then updates
$\lambda_{i+1} = f(x^{(i)}) / g(x^{(i)})$. The sequence of estimates
is strictly increasing and converges to $\lambda^*$ in finitely many
steps; for linear fractional combinatorial problems, the number of
iterations is strongly polynomial~\citep{radzik1998fractional}.

Beyond parametrization, exact methods include B\&B and cutting-plane
algorithms tailored to the fractional structure. Early B\&B approaches
exploited pseudo-Boolean or nondecreasing function
properties~\citep{robillard1971hyperbolic, agrawal77, saipe1975}, while
others used continuous relaxations to derive
bounds~\citep{bitranmag76}. Cutting-plane methods iteratively solve the
linear relaxation and add Gomory-style cuts to enforce integrality,
leveraging the fact that optima lie at extreme
points~\citep{granot1977integer, grunspan1973hyperbolic, Nemhauser99}.
More recent contributions include specialized branch-and-cut
frameworks~\citep{MENDEZDIAZ2014246}, Lagrangian
relaxation~\citep{Amaldi2012}, iterative reduction to sequences of
integer linear subproblems~\citep{anzai1974integer}, and
reformulations into equivalent MILPs solvable by commercial
solvers~\citep{tawarmalani2002global, Busygin2005, IBARAKI197639}.

Recent advances include denominator-objective
restriction~\citep{ponnaiah2013solving}, convexification
techniques~\citep{he2024convexification}, and novel relaxations for
quadratic binary and nonlinear mixed-integer fractional
problems~\citep{ADAMS200555, ADAMS2007510, adams86, boros2002pseudo,
Sherali2007, CHAOVALITWONGSE2004517, gupte13,
sherali2013reformulation}. In addition to exact methods, a variety of heuristic approaches have been proposed, including simulated annealing and tabu search~\citep{Hansen199097}, greedy randomized adaptive search procedures (GRASP)~\citep{prokopyev2005multiple}, and application-specific Lagrangian-based heuristics~\citep{amiri1999bandwidth}.


This chapter has reviewed three interconnected foundations for the dissertation. First, it demonstrated the importance of EHRs in analyzing comorbidity patterns and mortality risk. Second, it described how the maximum clique problem, together with its rich repertoire of enumerative, exact, and heuristic solution methods, provides a combinatorial tool for identifying densely interconnected groups of diseases. Third, it showed how fractional programming, particularly its extensions to integer and combinatorial settings, offers a mathematically rigorous framework for optimizing ratio-based objectives such as mortality rate of a disease cluster. Building on these foundations, this dissertation makes the following contributions.

\section{Contributions}\label{sec.contributions}

In Chapter~\ref{sec.probstmnt}, we introduce the maximum mortality
rate clique problem, a fractional combinatorial optimization
problem defined on a comorbidity graph derived from EHR data. This problem seeks an $\ell$-frequent clique of diseases that maximizes the mortality rate of its shared patient population.

In Chapter~\ref{sec.theory}, we prove that the mortality rate function
is neither additive nor monotonic, ruling out simple greedy strategies
and motivating the need for exact methods. We then establish
NP-completeness of the associated decision problem under several
parameterizations of the mortality threshold and clique size, and show
that even the simpler problem of finding any large $\ell$-frequent
disease subset, without requiring clique structure or a mortality
target, is NP-complete.

In Chapter~\ref{sec.methods}, we develop six solution methods of
increasing sophistication. We begin with a modified Bron--Kerbosch
enumerative algorithm and a combinatorial branch-and-bound
algorithm, the latter augmented with an upper-bounding scheme
that exploits the structure of the mortality rate function to prune
the search tree substantially. We then present two mixed-integer linear programming formulations:
one based on auxiliary product variable linearization and one based on
the Charnes--Cooper transformation. Additionally, we introduce a logic-based Benders decomposition that embeds naturally in a decomposition
branch-and-cut framework. We further develop a hybrid variant that
integrates the combinatorial branch-and-bound algorithm to generate valid inequalities, yielding substantially faster convergence and the ability to solve larger instances to optimality.

In Chapter~\ref{sec.DATAPREP}, we describe how the comorbidity graph and the experimental testbed are derived from real-world EHR data. This establishes the empirical foundation for the computational study, while the detailed extraction, code-mapping, and graph-construction procedures are presented in that chapter.

In Chapter~\ref{sec.results}, we evaluate all six methods on a testbed of instances derived from the MIMIC-IV EHR database. We demonstrate that the hybrid logic-based Benders decomposition algorithm solves large-scale instances to optimality while accommodating side constraints without algorithmic modification, illustrated through  case studies. We also introduce a size-restricted variant of the problem and a marginal mortality rate extension, and test both on the full MIMIC-IV cohort of \numprint{223291} patients, identifying clinically interpretable high-mortality disease cliques across a range of clique sizes.


Chapter~\ref{sec.conclusion} concludes the dissertation by summarizing
the contributions, discussing limitations, and outlining directions
for future research.

